\section{Nonzero-Gap Finite-Enclosure Optimization}
\label{sec:nonzero-gap-finite-enclosure-optimization}

This appendix proves the selected-gap and supported-rescuer enclosure
statements and the two-vertex replacement. The witnesses are fixed in
Section~\ref{sec:finite-enclosure}; the terminal enclosure proofs use hull-edge calipers rather than signed center parameters.


\subsection{Direct neighboring-ray domination}
\label{subsec:new-neighbor-ray-capacity}

The capacities keep their variational definitions. The following proof
establishes the bound needed by all active witnesses without evaluating
the cubic plateau and radical branches of the full function.


\begin{lemma}[Sharp diagonal neighboring capacity]
\label{lem:direct-neighbor-diagonal}
For $0\le m\le1/2$, one has $C_+(m,m)=C_-(m,m)=1-m$.
\end{lemma}
\begin{proof}
Use oblique coordinates with metric $x^2+y^2-xy$. The four points are
$P=(0,0)$, $A=(m,0)$, $B=(0,m)$, and $Z=(c,1)$. The case $m=0$
is immediate. If $0<m\le1/2$ and $c>1-m$, their hull has cyclic order
$P,A,Z,B$. \zcref{lem:support-cell-rotation} reduces the enclosure to its four outward
edge normals. At $PA$ and $BP$, the side lengths are $c+m$ and $1+m$,
both greater than one.

Embed $(x,y)$ as $(x-y/2,hy)$. At the other two edges take unnormalized
normals $n_A=(h,1/2-c+m)$ and $n_B=(h(m-1),m/2+c-1/2)$. Put
\[
\begin{aligned}
 D_A&=(c-m)^2-(c-m)+1,&N_A&=m+c^2-cm-c+1,\\
 D_B&=c^2+cm-c+m^2-2m+1,&N_B&=cm+c^2-c-m+1.
\end{aligned}
\]
Selecting the contained points $(A,Z,P)$ in the three normal directions
at $n_A$, and $(B,P,Z)$ at $n_B$, gives support lower bounds
$L_{AZ}\ge N_A/\sqrt{D_A}$ and $L_{ZB}\ge N_B/\sqrt{D_B}$.
No active-support selector is assumed. With $q=c+m-1>0$,
\[
\begin{aligned}
 N_A^2-D_A={}&q^2\big((q+1-3m)^2+4m^2-2m+1\big)\\
 &+q(1-2m)(6m^2-2m+1)+m^2(1-2m)^2>0,\\
 N_B^2-D_B={}&q^4+(2-2m)q^3+(m^2-4m+2)q^2\\
 &+(1-m)(1-2m)q>0.
\end{aligned}
\]
Here $4m^2-2m+1\ge3/4$, $6m^2-2m+1>0$, and
$m^2-4m+2\ge1/4$. The numerators are positive:
$N_A\ge(1+m)(3-m)/4$ and $N_B=1-m+(1-m)q+q^2$.
Thus every caliper requires side greater than one.
Equality is realized by the unit triangle
$\conv\{(-m,0),(1-m,0),(1-m,1)\}$. Reflection gives $C_-$.
\end{proof}




% Active elementary proof from numbered source 2615.
\subsection{Elementary local capacity bounds}
Write $f(a,b)=1-c_{\max}(a,b)$ for $a,b\ge0$, $a+b\le1$, and put
$m=\min(a,b)$. Let $h=\sqrt3/2$ and let $c(m)\in[h,1]$ be the root of
\[
 F(C,m)=C^4-C^2+mC-m^2=0,\qquad \ell(m)=1-c(m).
\]
It is unique because $F_C>0$ on $[h,1]$,
$F(h,m)=-(m-h/2)^2\le0$, and $F(1,m)=m(1-m)\ge0$.
Define
\[
 q(m)=\frac{m(1-m)}2,\qquad
 \beta(m)=\frac{2m}{4+3m},\qquad \kappa(m)=1+\frac32m.
\]
\begin{lemma}[Own-ray deficit bounds]
\label{lem:bc-capacity-tools}
For $a,b\ge0$, $a+b\le1$, and $m=\min(a,b)$,
\[
 q(m)\le\ell(m)\le f(a,b)\le m.
\]
For $m\le1/4$, also $\beta(m)\le\ell(m)$ and $\beta(m)\ge q(m)$.
\end{lemma}
\begin{proof}
The upper bound and capacity antitonicity follow from
\zcref{lem:new-common-pair-domination}. By the selected local formula
\zcref{prop:new-exact-local-set}, for $m\le h/2$ the diagonal pair lies
in the quartic cell (its cell polynomial is $m^2(16m^2-3)\le0$), so
$c_{\max}(a,b)\le c_{\max}(m,m)=c(m)$. For $m\ge h/2$,
antitonicity gives $c_{\max}(a,b)\le c_{\max}(h/2,h/2)=h\le c(m)$.
Thus $f(a,b)\ge\ell(m)$.

For $c=c(m)$, the root equation and a factorization give
\[
 \ell c^2(1+c)=m(c-m),
\]
\[
 2(c-m)-(1-m)c^2(1+c)
 =(1-c)\bigl((1-m)c(c+2)-2m\bigr)\ge0.
\]
The bracket is at least
$\tfrac12\cdot\tfrac34\cdot\tfrac{11}4-1=1/32$, since $c\ge h>3/4$
and $m\le1/2$. Division by $c^2(1+c)>0$ proves $\ell\ge q(m)$.

For the stronger bound on $m\le1/4$, put
\[
 B(m)=81m^4+405m^3+656m^2+272m-128.
\]
Direct substitution gives
\[
 F(1-\beta(m),m)=-\frac{m^2B(m)}{(4+3m)^4}.
\]
The polynomial $B$ is increasing and $B(1/4)=-3163/256<0$.
Also $1-\beta(m)\ge17/19>h$. Monotonicity of $F$ in $C$ implies
$\beta(m)\le\ell(m)$. Finally
$\beta(m)-q(m)=m^2(3m+1)/(2(3m+4))\ge0$. At $m=0$ all bounds are immediate.
\end{proof}

\begin{lemma}[Quarter-range slack envelope]
\label{lem:bc-slack-envelope}
For $a,b\ge0$, $a+b\le1$, $m=\min(a,b)\le1/4$, and $\delta=1-a-b$,
\[
 f(a,b)\ge\max\{\beta(m),\ m-\kappa(m)\delta\}.
\]
The two entries meet at $\delta=\beta(m)$. This estimate is asserted only
on $m\le1/4$; the broader historical envelope is not needed here.
\end{lemma}
\begin{proof}
Since $m-\kappa(m)\beta(m)=\beta(m)$, the baseline in
\zcref{lem:bc-capacity-tools} proves the claim for $\delta\ge\beta(m)$.
The case $m=0$ is immediate. Suppose $m>0$ and $0<\delta\le\beta(m)$.
Put $M=1-m-\delta$. Then $M\ge m$ (indeed
$1-2m-\beta(m)\ge15/38>0$) and $1-\delta\ge1-\beta(m)\ge c(m)$.
The selected triangular branch of \zcref{prop:new-exact-local-set}
therefore applies: the capacity is the smaller root of
\[
 H(C)=\bigl((1-\delta)^2-1\bigr)C^2+MC-M^2.
\]
At $\overline C=1-m+\kappa(m)\delta$, expansion gives
$H(\overline C)=\delta N_m(\delta)/4$, where
\[
\begin{split}
 N_m(\delta)={}&(3m+2)^2\delta^3-10m(3m+2)\delta^2\\
 &+(28m^2-22m-20)\delta+14m(1-m).
\end{split}
\]
Because $\delta\le\beta(m)\le m/2\le1/8$,
\[
 N_m'(\delta)\le3(11/4)^2(1/8)^2+7/4-20=-18325/1024<0.
\]
At the right endpoint,
\[
 N_m(\beta(m))=-\frac{2mB(m)}{(4+3m)^3}>0.
\]
Thus $H(\overline C)>0$. The quadratic is concave and the capacity is
its selected smaller root, so $c_{\max}(M,m)<\overline C$. This proves
the affine lower bound. At $\delta=0$, the exact capacity is $1-m$.
\end{proof}

\begin{lemma}[Coupled nonuniform capacity inequality]
\label{lem:bc-coupled-capacity}
If $0<x\le y<1$ and $x/2\le z\le y$, put
\[
 r=f(1-y,z),\qquad t=f(1-z,x/2).
\]
Then
\[
 2(r+t)>(1-y+r)(2x-y+r).
\]
\end{lemma}
\begin{proof}
Put $u=x/2$ and
\[
 P(y,r,t)=2(r+t)-(1-y+r)(4u-y+r).
\]
The domain is $0<u$, $2u\le y<1$, $u\le z\le y$ and
$0<r,t\le1/2$. By \zcref{lem:bc-capacity-tools},
\[
 P_r=1-4u+2y-2r\ge1-2r\ge0,\qquad P_t=2.
\]
Hence valid lower bounds on either radius may be substituted. These are
scalar comparisons; the exact witness radii are not changed.

\readerheading{Easy regions.}
First let $y\ge1/2$ and put $a=1-y$. Then $0<a\le1/2$,
$u\le(1-a)/2$, and $u\le z\le1-a$.
If $z\ge a$ and $z\le1-u$, substitute $q(a),q(u)$.
The expression is concave in $u$; at $u=0$ it equals
$2q(a)+(a+q(a))(1-a-q(a))>0$, and at $u=(1-a)/2$ it equals
$(1-a)^3(1+a)/4>0$.
If $z\ge a$ and $z\ge1-u$, both radii are at least $q(a)$ because
$a\le1-z\le u\le1/2$ and $q$ is increasing. The substituted expression
decreases with $u$ and at $u=(1-a)/2$ equals
$a(1-a)(a^2-a+2)/4>0$.
If $z\le a$, both radii are at least $q(u)$. The expression is concave
in $a\in[u,\min(1/2,1-2u)]$, with endpoint values
\[
\begin{array}{c|c|c}
a&P\text{ after substitution}&u\text{ range}\\\hline
u&u(14-43u+14u^2-u^3)/4&(0,1/3]\\
1/2&(1-17u^2+10u^3-u^4)/4&(0,1/4]\\
1-2u&u(-u^3+2u^2+9u-2)/4&[1/4,1/3].
\end{array}
\]
The three brackets respectively decrease to $32/27$, decrease to
$23/256$, and increase from $23/64$, so all are positive.

Now let $y\le1/2$ and $z\ge2u$. Substitute $q(z),q(u)$.
The result increases with $y$, so set $y=z$; it is then concave in
$u\in[0,z/2]$. The value at zero is positive, and the value at
$u=z/2$ is $z^3(2-z)/4>0$.
Finally, for $y\le1/2$ and $1/4\le z\le2u$, the same substitution
allows $y=2u$, giving
\[
 P\ge q(z)(1-q(z))-u+3u^2
 \ge\frac3{32}\frac{29}{32}-\frac1{12}=\frac5{3072}>0.
\]
Here $q(z)(1-q(z))$ increases on $[1/4,1/2]$ and
$-u+3u^2=3(u-1/6)^2-1/12$.

\readerheading{The remaining region and its switches.}
It remains to treat
\[
 0<u\le1/4,\quad u\le z\le\min(2u,1/4),\quad 2u\le y\le1/2.
\]
The smaller capacity coordinates are $z$ and $u$. By
\zcref{lem:bc-slack-envelope} compare with
\[
 R_0=\max\{\beta(z),z-\kappa(z)(y-z)\},\qquad
 T_0=\max\{\beta(u),u-\kappa(u)(z-u)\}.
\]
Put $g(v)=v+\beta(v)$, an increasing rational function. The switches are
$y=g(z)$ and $z=g(u)$.
For fixed $u,z$, the affine branch of $R_0$ gives
\[
 \frac{d^2}{dy^2}P(y,R_0,T_0)=-2(1+\kappa(z))^2<0.
\]
The baseline branch gives $dP/dy=1+4u+2\beta(z)-2y>0$.
Since $g(z)\le g(1/4)=27/76<1/2$, only $y=2u$ or $y=g(z)$
needs to be checked, whenever the latter belongs to the domain.

On $y=2u$,
\[
 P=R_0(1-R_0)+2T_0-2u(1-2u).
\]
Split the $z$ interval at $g(z)=2u$ and $z=g(u)$.
The function $\beta(z)(1-\beta(z))$ is concave, since $\beta$ is concave
and $v(1-v)$ is increasing and concave on $[0,1/2]$.
On the other branch,
\[
 R=2z+\tfrac32z^2-2u-3uz,\qquad R'=2+3z-3u\ge2,\qquad R''=3.
\]
Where it is active, $0<R\le z\le1/4$, so
\[
 (R(1-R))''=3(1-2R)-2(R')^2\le-5<0.
\]
As $T_0$ is affine or constant on each piece, $P$ is piecewise concave.
At $z=u$ it equals $\beta(u)(1-\beta(u))+4u^2>0$; at $z=2u$ it
is $2T_0>0$. The endpoint $z=1/4$ is covered by the easy-region bound,
since the comparison radii dominate the $q$ bounds. Only the switches remain.

On $y=g(z)$,
\[
 P=2\beta(z)+2T_0-(1-z)(4u-z),\qquad z/2\le u\le g(z)/2.
\]
Its $u$ derivative on the baseline branch is
$16/(4+3u)^2-4(1-z)\le-2<0$; on the other branch it is $6u+z>0$.
Thus the minimum lies at an endpoint or $z=g(u)$. The endpoints are
$z=2u$, already checked, and transition A below. The ordering of switches
is not assumed; coincident switches are included.

\readerheading{Three transition checks.}
In A, $y=2u=g(z)$. Using the affine entry as a lower bound for $T_0$ gives
\[
 P\ge\frac{z^2(9z^2+36z+44)}{4(3z+4)^2}>0.
\]
In B, $y=2u$ and $z=g(u)$, so $T_0=\beta(u)$. If $R_B$ denotes the
affine entry of $R_0$, then
\[
 R_B-u(1-2u)=\frac{u^2(9u^2+6u+4)}{2(3u+4)^2}>0.
\]
Since $v(1-v)$ increases on $[0,1/2]$,
\[
 P\ge\frac{u^2(1+19u-4u^2-12u^3)}{3u+4}>0,
\]
where $4u^2+12u^3\le7/16$.
In C, $y=g(z)$ and $z=g(u)$, direct substitution gives
\[
 P=\frac{3u^2(81u^4+441u^3+768u^2+460u+48)}
 {(3u+2)(3u+4)^2(3u+8)}>0.
\]
These regions and transitions are exhaustive. Monotonicity in the radii
therefore proves the inequality for the exact capacities.
\end{proof}

% Geometry independent of the capacity construction.
\subsection{A five-point threshold criterion}
\begin{lemma}[Five-point threshold criterion]
\label{lem:bc-five-point}
\label{lem:five-point-threshold}
Suppose $0<x\le y<1$, $0<r\le\min(y,1-y)$, and $0<t\le x/2$.
Put $G_w=(1-w)V_0+wV_1$, $k=y-r$, and
$F=\{M_0,G_x,G_y,rV_2,tV_4\}$. Then
\begin{equation}
 \Lambda(F)>1\quad\Longleftrightarrow\quad
 r+t+\max\{1/2,1-x(1-k)\}>\sqrt{1-k+k^2}.
 \label{eq:five-point-threshold}
\end{equation}
The gap points may coincide. This is a criterion at side one, not a
claim that the displayed caliper always minimizes the enclosing side.
\end{lemma}
\begin{proof}
Put $M=M_0$, $P=rV_2$, and $Q=tV_4$. For $x<y$, these points have
hull order $M,G_x,G_y,P,Q$. The stated bounds make the consecutive
cross products positive. At $G_yP$, the outward normal $n=(-hk,1-k/2)$ has
squared norm $1-k+k^2$. The three supports divided by $h$ are
$r$, $t$, and $\max\{1/2,1-x(1-k)\}$. Hence its caliper side is
\[
 L_3=\frac{r+t+\max\{1/2,1-x(1-k)\}}{\sqrt{1-k+k^2}}.
\]
The other four calipers exceed one. At the edges $MG_x$, $G_xG_y$,
and $QM$, selected point projections give
\[
 L_1>\frac{1-x+2x^2}{\sqrt{1-2x+4x^2}},\qquad
 L_2=1+r+t,\qquad
 L_5\ge\frac{1+t+2t^2}{\sqrt{1+2t+4t^2}}.
\]
For $L_1$ use $y\ge x$ and the strictly positive contribution $2tx$;
for $L_5$ use $y\ge x\ge2t$. The squared differences are
\[
 x^2(2x-1)^2\ge0,\qquad t^2(1+2t)^2>0,
\]
respectively. At $PQ$, put $v=t/r>0$. Selected projections give
\[
 L_4\ge\frac{t+v+x+\max\{1/2,1-(1+v)x\}}{\sqrt{1+v+v^2}}.
\]
The piecewise-linear expression in $x$ is minimized at $1/[2(1+v)]$.
Thus its numerator exceeds $g(v)=v+1/2+1/[2(1+v)]$, and
\[
 g(v)^2-(1+v+v^2)=\frac{v^2}{4(1+v)^2}>0.
\]
By \zcref{thm:cert-caliper}, only $L_3$ can cross the unit threshold.
When $x=y$, discard the zero edge $G_xG_y$; the other four calculations
remain valid, without a limiting strictness argument.
\end{proof}

\readerheading{Capacity instance.}
For $x/2\le z\le y$, take $r=f(1-y,z)$ and $t=f(1-z,x/2)$.
The bounds of \zcref{lem:bc-capacity-tools} meet the geometric hypotheses,
and \zcref{lem:bc-coupled-capacity} gives $r+t>(1-k)(x-k/2)$.
Consequently the left side of \eqref{eq:five-point-threshold} exceeds
$1-k/2+k^2/2\ge\sqrt{1-k+k^2}$; the squared difference is
$k^2(1-k)^2/4$. Thus this capacity instance has $\Lambda(F)>1$.



\subsection{The supported-rescuer adapter}



\begin{lemma}[Supported endpoint to four-point enclosure]
\label{lem:new-rescuer-tail-budget}
Suppose $T_0$ has actual boundary endpoint $a=A_0$ toward $V_5$ and
supported interval $[c,u]$ on $r_1$, measured from $V_1$ toward $O$.
Put $\varepsilon=1-u>0$ and $M=M_c^{\rm sup}$. Assume
\[
 \varepsilon V_1\in U_C,\qquad A_0+\varepsilon\le1,
 \qquad \frac{A_0}{A_0+\varepsilon}\le1-M.
\]
If $T_1$ is uniquely supercritical with $C_1\ge c$, and the
four center-free edges from $e_{1,2}$ through $e_{4,5}$ join nonsupercritical
$T_2,T_3,T_4,T_5$, then the skeleton is not covered. The conclusion applies
to both possible nonzero gap ranks.
\end{lemma}
\begin{proof}
The strict-supercritical envelope gives $B_1<M$, and the four handoffs give
$B_5\le B_4\le B_3\le B_2\le B_1$. Thus $B_5<1-A_0$, so the left-edge
actual gap has endpoints $Y(A_0),Y(1-B_5)$, with $Y(t)=(1-t)V_0+tV_5$.
Under coverage these endpoints and $O,\varepsilon V_1$ belong to $U_C$.
Since $M\ge1/2$, the hypotheses imply
\[
 A_0\le\varepsilon,\qquad
 B_5<M\le\frac{\varepsilon}{A_0+\varepsilon}.
\]
Apply \zcref{thm:fixed-four-point-rescuer} with geometric parameters
$a=A_0$, $\beta=B_5$. The fixed set has enclosure number at least one,
contrary to compact containment in $U_C$.
\end{proof}

\subsection{T3-like supported-tail adapter}

\begin{lemma}[Solved T3-like supported-tail adapter]
\label{lem:appendix-t3-supported-tail}
Assume a nonzero-gap skeleton cover has center midpoint $M_0$, a T3-like
role $T_0$ supplying $M_1$, uniquely supercritical $T_1$, and nonsupercritical
Vd0 roles at the other four vertices. The actual supported interval
satisfies the hypotheses of \zcref{lem:new-rescuer-tail-budget}.
\end{lemma}
\begin{proof}
Use the original Type-II chart \eqref{eq:orientation-type-II}, with
positive first slack $\alpha_T$ and actual reach $a=A_0$.
Thus $0<t<1$, $z=\sqrt{1-t+t^2}$, and substitution on $r_1$ gives
\begin{equation}
 c=\frac{1+\alpha_T+a-z}{1-t},\qquad
 u=a+t,\qquad \varepsilon=1-u.
 \label{eq:t3-actual-supported-interval}
\end{equation}
Open midpoint containment gives $0<c<1/2<u<1$.
The supercritical own role ends before the midpoint and the other roles
have no positive trace on $r_1$. Thus $u$ is the total radial frontier:
$\varepsilon V_1\in U_C$ by \zcref{lem:fixed-total-radial-forcing}.
Coverage of the near endpoint gives $C_1\ge c$, since the center exits
before $M_1$.

Put $x=a/(1-t)$ and $\theta=t/(1+z)$. Then
\begin{equation}
 a+\varepsilon=1-t<1,\qquad
 c=x+\theta+\frac{\alpha_T}{1-t}>x+\theta,\qquad
 \frac a{a+\varepsilon}=x.
 \label{eq:t3-actual-ratio}
\end{equation}
In particular $0<x<c<1/2$. Since $z<1$ and
$x+(1-x)t=a+t>1/2$,
\[
 \theta>\frac t2>\frac{1-2x}{4(1-x)}.
\]
The quadratic $Q_c(x)=x^2+(c-2)x+c$ increases with $c$, so
\[
 Q_c(x)>2x^2+(\theta-1)x+\theta
 >\frac{(1-2x)(4x^2-3x+1)}{4(1-x)}>0.
\]
Here $4x^2-3x+1=4(x-3/8)^2+7/16$.
The scalar rescuer test \zcref{lem:rescuer-ratio-test} gives $x\le1-M_c^{\rm sup}$.
Together with the actual size relation and the forced endpoint, these
are exactly the common D hypotheses. No translated near endpoint is used.
\end{proof}

\begin{lemma}[Solved Vd1 supported-tail adapter]
\label{lem:appendix-vd1-supported-tail}
In the center-based Vd1 supplier placement, with $M_1\in U_0$ and
$T_1$ uniquely supercritical, the actual corner data imply the size and
ratio inequalities required by \zcref{lem:new-rescuer-tail-budget}.
\end{lemma}
\begin{proof}
Set $a=A_0$, $b=B_0$. By \zcref{lem:vd1-forward-orientation}, the
original corner parameter satisfies $t\ge1$. Put
$d=\sqrt{t^2+t+1}$ and write the actual supported interval as $[c,u]$.
Then $0<c<1/2<u<1$ and
\[
 c=\frac{t(1-b)}{t+1},\qquad
 u=\frac{d-a-tb-1}{t}.
\]
The midpoint condition and $tb=t-c(t+1)$ imply
\[
 a\le F(t,c):=d-1-\frac{3t}{2}+c(t+1).
\]
Since $d'<1$ and $c<1/2$, $F$ decreases with $t$, so
\[
 a\le F(t,c)\le F(1,c)=\sqrt3-\frac52+2c=:L(c).
\]
Put $\varepsilon=1-u$. Direct substitution gives
\[
 \varepsilon=\varepsilon_0+\frac at,\qquad
 \varepsilon_0=\frac{2t+1-c(t+1)-d}{t}>0.
\]
For the last sign, $d<t+1$ gives
$t\varepsilon_0>(1-c)t-c\ge1-2c>0$.
The function $v/(v+\varepsilon_0+v/t)$ increases for $v\ge0$, and
$\varepsilon_0+F(t,c)/t=1/2$. Hence, with $x=a/(a+\varepsilon)$,
\[
 x\le\frac{F(t,c)}{F(t,c)+1/2}\le2F(t,c)\le2L(c)
 <4c-\frac32,
\]
using $\sqrt3<7/4$. Thus $0<x<1/2$ and $c>(3+2x)/8$. Therefore
\[
 x^2+(c-2)x+c>
 x^2+\left(\frac{3+2x}{8}-2\right)x+\frac{3+2x}{8}
 =\frac{(1-2x)(3-5x)}8>0.
\]
The scalar rescuer test \zcref{lem:rescuer-ratio-test} proves $x\le1-M_c^{\rm sup}$.

Positive forward support also gives
\[
 (t+1)a+tb<(t+1)(d-1)-t^2<d-1,
\]
because $d<t+1$. Thus $u>a$ and $a+\varepsilon=1+a-u<1$.
As in the T3-like placement, $u$ is the total frontier and the near
endpoint is supplied by $T_1$, so $\varepsilon V_1\in U_C$ and
$C_1\ge c$. These are all the common D hypotheses.
\end{proof}

\subsection{Vd1 two-chart replacement}

The replacement below uses two different vertex charts.  It preserves the
five skeleton pieces that can change, then invokes the all-nonsupercritical
obstruction without assuming preservation of the input gap rank.

\begin{lemma}[Half-square admissibility]
\label{lem:appendix-half-square-admissibility}
If \((x,y,z)\in\mathcal A\) satisfies
\[
 x\ge\frac12,
 \qquad
 0\le z\le\frac12,
\]
then
\[
 y\le1-z.
\]
\end{lemma}

\begin{proof}
Suppose \(y>1-z\).  Then \(x+y>1\), so the selected supercritical cell of
\zcref{prop:new-exact-local-set} applies.  In the ordered half \(x\le y\)
its necessary polynomial is
\[
 F(x,y,z)=
 (x^2-1)z^2+(2xy^2+y)z+y^4-y^2\le0.
\]
It is nondecreasing in \(x\), its derivative in \(y\) is positive for
\(y\ge1-z\), and
\[
 F\left(\frac12,1-z,z\right)
 =\frac{z^2(2z-5)(2z-1)}4\ge0.
\]
This is a contradiction.  In the reflected half \(y<x\), apply the same
argument to \(F(y,x,z)\).  At the joining corner,
\[
 F(1-z,1-z,z)=z(1-2z)\ge0,
\]
so no selected feasible point is lost at equality.
\end{proof}

\begin{lemma}[Two-vertex replacement from scalar margins]
\label{lem:two-vertex-replacement}
Consider two adjacent V roles. Only the five pieces
\[
 e_{5,0},\quad e_{0,1},\quad e_{1,2},\quad r_0,\quad r_1
\]
may be affected by their replacement, and the shared edge is center-free.
Let $a,c,B,r\in[0,1)$ be the boundary reach on $e_{5,0}$, the own-radial
reach on $r_0$, the boundary reach on $e_{1,2}$, and the radial demand
on $r_1$ needed to overlap the unchanged center interval. If
\[
 a<\tfrac12,\qquad a+c<1,\qquad a+B<1,\qquad
 r<\max\{1-a,1-B\},
\]
there are two open unit equilateral replacements preserving these five
skeleton pieces, with actual boundary sums below one and no positive
adjacent support. No center-class formula is assumed.
\end{lemma}
\begin{proof}
Choose first $p_2\in(a,1-B)$ with $\max\{p_2,1-p_2\}>r$. This is
possible because the supremum on that interval is $\max\{1-a,1-B\}$.
Next choose $a<p_1<\min\{p_2,1/2,1-c\}$ and
\[
 0<\varepsilon<\min\{p_1-a,p_2-p_1,1-B-p_2,
                         1-p_1-c,\max(p_2,1-p_2)-r\}.
\]
All five margins are positive. Use the separate vertex charts
\[
 X_0(x,y)=V_0+x(V_5-V_0)+y(V_1-V_0),\qquad
 X_1(x,y)=V_1+x(V_0-V_1)+y(V_2-V_1).
\]
Both have metric $x^2+y^2-xy$. Define
\[
 \Delta_p^-=\conv\{(0,1-p),(1,1-p),(0,-p)\},\quad
 D^-_{p,\varepsilon}=\interior(\Delta_p^-)+(-\varepsilon,0)
 \quad(p\le1/2),
\]
\[
 \Delta_p^+=\conv\{(p,0),(p,1),(p-1,0)\},\quad
 D^+_{p,\varepsilon}=\interior(\Delta_p^+)+(0,-\varepsilon)
 \quad(p>1/2).
\]
Their edges have unit length. The minus half-planes are
$x>-\varepsilon$, $y<1-p$, $y>x+\varepsilon-p$; the plus half-planes
are $y>-\varepsilon$, $x<p$, $x>y+\varepsilon+p-1$.
Substitution on the axes and diagonal gives their actual reaches:
\[
\begin{array}{c|ccc}
 &A&B&C\\\hline
 D^-_{p,\varepsilon}&p-\varepsilon&1-p&1-p\\
 D^+_{p,\varepsilon}&p&1-p-\varepsilon&p
\end{array}
\]
Both contain the chart origin. On the adjacent support lines $(z,1)$ and
$(1,z)$ the half-planes are incompatible for $0\le z\le1$, so neither
has positive adjacent support.
Set $U'_0=X_0(D^-_{p_1,\varepsilon})$ and use
$U'_1=X_1(D^-_{p_2,\varepsilon})$ for $p_2\le1/2$, or
$U'_1=X_1(D^+_{p_2,\varepsilon})$ otherwise. The two charts and this
split cannot be identified. Both actual boundary sums are $1-\varepsilon$.

The first replacement reaches $p_1-\varepsilon>a$ on $e_{5,0}$ and
$1-p_1>c$ on $r_0$. On the shared edge the reaches sum to at least
\[
 (1-p_1)+(p_2-\varepsilon)>1,
\]
so the open traces overlap. The second reaches more than $B$ on
$e_{1,2}$ and $\max(p_2,1-p_2)>r$ on $r_1$. Thus all five pieces
remain covered; all other pieces are unchanged.
\end{proof}

\begin{proposition}[Vd1 input to the two-vertex replacement]
\label{prop:appendix-vd1-two-chart-replacement}
Assume a skeleton cover has a unique Vd1 role $T_0$ supplying $M_1$, a
unique supercritical role $T_1$, and all other V roles are nonsupercritical
Vd0. Suppose $e_{0,1}$ is center-free and no other V role has positive
trace on $r_0$ or $r_1$. The scalar replacement applies in CE1 and CE2
and yields a skeleton cover with $N'_+=0$.
\end{proposition}

\begin{proof}
Let \(A_0,B_0,C_0\) be the actual incoming, outgoing, and own-radial
reaches of the Vd1 role.  Its corner normal form and \zcref{lem:vd1-forward-orientation} supply
\(\vartheta\ge1\) and
\[
 \omega:=\sqrt{\vartheta^2+\vartheta+1}
\]
such that
\[
 C_0=\frac{\omega-A_0-\vartheta B_0}{\vartheta+1},
 \qquad
 \lambda=\frac{\vartheta(1-B_0)}{\vartheta+1},
\]
\[
 u_{\rm adj}
 =\frac{\omega-A_0-\vartheta B_0-1}{\vartheta}.
\]
The strict midpoint and one-sided-support conditions give
\[
 A_0+C_0<1,
 \qquad
 A_0<\lambda\le\frac12,
\]
\[
 B_0<\frac12,
 \qquad
 u_{\rm adj}<1-A_0.
\]
For the first inequality, evaluation on the closed midpoint face reduces
strictness to
\[
 (2\vartheta^2+4\vartheta+1)^2
 -4(\vartheta+1)^2(\vartheta^2+\vartheta+1)
 =4\vartheta^3+4\vartheta^2-4\vartheta-3>0.
\]
The remaining inequalities follow directly from the displayed endpoint
formulas and the strict midpoint inequalities.

The center-free shared edge and the supported interval give
\[
 A_1\ge1-B_0>\frac12,
 \qquad
 C_1\ge\lambda
\]
for the actual reaches of \(T_1\).  The selected supercritical component
has \(C_1\le1/2\), so
\zcref{lem:appendix-half-square-admissibility} yields
\[
 B_1\le1-C_1\le1-\lambda.
\]
Since \(A_0<\lambda\),
\[
 A_0+B_1<1.
\]

Define the C-triangle reach on \(r_1\), measured from \(O\), by
\[
 d_1^C:=\max\{x\in[0,1]:xV_1\in T_C\},
 \qquad
 c_1^{\rm req}:=1-d_1^C.
\]
Only \(U_C,U_1\), and the adjacent trace of \(U_0\) can contribute a
positive interval to \(r_1\).  Open skeleton coverage gives
\[
 c_1^{\rm req}<\max\{C_1,u_{\rm adj}\}
 \le\max\{1-B_1,1-A_0\}.
\]

Apply \zcref{lem:two-vertex-replacement} with
\[
 (a,c,B,r)=(A_0,C_0,B_1,c_1^{\rm req}).
\]
All four scalar conditions were proved above. The four untouched roles
remain nonsupercritical Vd0, while the two new actual boundary sums are
$1-\varepsilon<1$. The resulting skeleton cover contradicts N0.
The original C triangle is fixed; its class and the input gap rank were
not used in the construction. This operation does not claim to preserve
coverage of the interior of $H$.
\end{proof}
